% ===================================================================
%  TAULA N.º 14 — El Carrer. — Els Comerciants.
%  Traduction 2026 d'apres texto/io, controlee sur texto/fr et
%  texto/es. Consigne : voir texto/ca/10-taula-01.tex.
%
%  LE TABLEAU LE PLUS CHARGE DU LIVRET : quatre-vingt-dix-neuf
%  renvois. Le suedois et le finnois y avaient chacun deux inversions
%  a defaire ; le catalan n'en a pas une, parce qu'il postpose son
%  complement du nom comme l'ido postpose le sien.
% ===================================================================

%%K t14-apar-1 apar
\VUsaut{4.0mm}
\VUtitre{15.0pt}{60}{TAULA N.º 14}
\VUsaut{3.0mm}

%%K t14-tit-1 sub
\VUpk{11.4pt}{40}{El Carrer. --- Els Comerciants.}
\VUsaut{3.0mm}

%%K t14-01-1 p
1. --- El meu amic Joan, que viu al camp, ha vingut a passar tres dies amb la
meva família. \VUgras{Fins} ara no havia tingut mai ocasió de veure una gran
ciutat, de manera que passem tots els dies passejant i jo li explico el que no
coneix.

%%K t14-02-1 p
2. --- Primer vam anar a veure l'\VUgras{observatori} \textsuperscript{(1)}, que li
va interessar moltíssim. En tornar pel \VUgras{carrer} \textsuperscript{(3)} on són
els \VUgras{banys turcs} \textsuperscript{(2)}, ens vam creuar amb un
\VUgras{enterrament} \textsuperscript{(4)}. Al capdavant del \VUgras{seguici
fúnebre} caminava el \VUgras{clergat}, davant del \VUgras{cotxe fúnebre} en què
s'havia col·locat el \VUgras{taüt}. Molts amics seguien la família de
\VUgras{dol}.

%%K t14-02-2 p
Quan hagué passat el seguici, vam creuar el carrer davant de la gran
\VUgras{cerveseria} \textsuperscript{(5)}, on molts \VUgras{clients} bevien
\VUgras{cervesa} a la terrassa, a l'abric de la \VUgras{marquesina}
\textsuperscript{(6)} envidrada.

%%K t14-03-1 p
3. --- A Joan el va intrigar l'\VUgras{escut d'armes} col·locat entre la botiga
del \VUgras{comerciant de vins i licors} \textsuperscript{(7)} i la del
\VUgras{venedor de música} \textsuperscript{(10)}. Em va preguntar què significava;
li vaig dir que assenyalava el \VUgras{consolat} \textsuperscript{(8)} anglès, i li
vaig mostrar el \VUgras{cònsol} \textsuperscript{(9)} dret al balcó.

%%K t14-tit-2 sub
\VUpk{10.2pt}{40}{Mirant els aparadors.}
\VUsaut{3.0mm}

%%K t14-04-1 p
4. --- Naturalment ens vam aturar davant l'exposició d'\VUgras{instruments de
música}, \VUgras{harmòniums}, \VUgras{orgues} i \VUgras{pianos}; vam mirar les
cobertes il·lustrades de les \VUgras{partitures} i les \VUgras{peces} per a piano,
i vam admirar la \VUgras{lira} de fusta daurada que servia de mostra.

%%K t14-05-1 p
5. --- Els núvols de pols que aixecaven els \VUgras{escombriaires}
\textsuperscript{(11)} i l'\VUgras{escombradora} \textsuperscript{(12)} ens van
obligar a passar de pressa per davant dels bells \VUgras{jocs de mobles} del
\VUgras{moblista} \textsuperscript{(13)} i de l'exposició d'\VUgras{estufes} i
\VUgras{làmpades} del \VUgras{ferreter} \textsuperscript{(14)}.

%%K t14-06-1 p
6. --- En aquell lloc, en efecte, vam haver de deixar la vorera, car estava
obstruïda pels embalums que es descarregaven d'un \VUgras{carro de mudances}
\textsuperscript{(16)} per portar-los a una \VUgras{botiga} \textsuperscript{(15)}
que acabava de llogar-se.

%%K t14-06-2 p
Això ens va impedir de veure els bells cotxes del \VUgras{carrosser}
\textsuperscript{(17)}, però no que ens aturéssim a mirar l'\VUgras{embalador}
\textsuperscript{(19)} que feia una caixa per al seu veí, el \VUgras{venedor de
bicicletes i automòbils} \textsuperscript{(18)}. Després, com que ja era prou tard,
vam tornar a casa.

%%K t14-tit-3 sub
\VUpk{10.2pt}{40}{Un passeig per la ciutat.}
\VUsaut{3.0mm}

%%K t14-07-1 p
7. --- L'endemà la meva mare va proposar que Joan es fes un retrat, i em va
demanar que l'acompanyés a casa del \VUgras{fotògraf} \textsuperscript{(20)}.
Després del dinar vam anar a l'\VUgras{estudi} de l'artista, on vam haver
d'esperar una bona estona. Per passar el temps vam mirar els \VUgras{retrats}
\textsuperscript{(22)} penjats a la paret.

%%K t14-07-2 p
Quan ens va arribar el torn la cosa no va durar gaire, i tan bon punt el meu amic
va acabar de posar davant de la \VUgras{càmera} \textsuperscript{(21)}, vam baixar.

%%K t14-08-1 p
8. --- Al primer pis vam veure, per la porta oberta de la \VUgras{perruqueria}
\textsuperscript{(23)}, l'oficial que tallava el cabell a un client.

%%K t14-08-2 p
Ja de nou al carrer, vam passar per davant de l'\VUgras{estanc}
\textsuperscript{(24)}. A Joan el van divertir tant el \VUgras{fanal}
\textsuperscript{(25)} vermell i la \VUgras{gran pipa} que servia de
\VUgras{mostra} \textsuperscript{(26)} que va ensopegar amb dos ancians que
xerraven prenent una \VUgras{pols de rapè} \textsuperscript{(27)}.

%%K t14-tit-4 sub
\VUpk{10.2pt}{40}{La pastisseria.}
\VUsaut{3.0mm}

%%K t14-09-1 p
9. --- Els \VUgras{bitllets} i les \VUgras{monedes} de tots els països exposats a
l'\VUgras{aparador} del \VUgras{canviador} \textsuperscript{(29)} ens van retenir un
moment, però com que era l'hora de berenar vam entrar a la \VUgras{pastisseria}
\textsuperscript{(31)} sense aturar-nos més.

%%K t14-10-1 p
10. --- Davant de totes les \VUgras{llaminadures} que hi havia a la botiga ---
\VUgras{pastissos, pa de pessic, galetes} i \VUgras{dolços} de tota mena --- amb
prou feines podíem triar. Al final Joan es va decidir pels \VUgras{pastissos de
crema}, i jo vaig prendre només una petita \VUgras{tarta de cireres}, car els
dolços \VUgras{em fan mal de queixal} i no tinc cap ganes de visitar massa sovint
el \VUgras{dentista} \textsuperscript{(30)}.

%%K t14-tit-5 sub
\VUpk{10.2pt}{40}{L'escriptura a màquina.}
\VUsaut{3.0mm}

%%K t14-11-1 p
11. --- Per ensenyar al meu amic una \VUgras{màquina d'escriure}, de la qual havia
sentit parlar però que no havia vist mai, vam entrar al \VUgras{despatx}
\textsuperscript{(32)} del meu \VUgras{cosí} \textsuperscript{(34)}. El vam trobar
dictant les seves cartes a la seva \VUgras{secretària} \textsuperscript{(35)}, que
és \VUgras{taquígrafa}. Amb prou feines acabada una carta, un
\VUgras{mecanògraf} \textsuperscript{(33)} la copiava a màquina. Com que no volíem
interrompre la feina del meu parent, ens hi vam quedar només uns minuts.

%%K t14-12-1 p
12. --- Sota el despatx del meu cosí viu un gran \VUgras{rellotger i joier}
\textsuperscript{(37)} que és a més argenter. La seva esplèndida botiga conté molts
\VUgras{rellotges de paret, rellotges de sobretaula, rellotges de butxaca},
\VUgras{cadenes} i \VUgras{joies} valuoses, com \VUgras{collarets de perles},
\VUgras{polseres} d'or, \VUgras{arracades} i \VUgras{anells} amb \VUgras{pedres
precioses} i \VUgras{diamants}. Un dels aparadors està dedicat del tot a
l'\VUgras{argenteria} i conté una gran col·lecció de \VUgras{coberts} de plata.

%%K t14-13-1 p
13. --- En doblar la cantonada del carrer, vaig notar que havien canviat la
\VUgras{placa} \textsuperscript{(36)} que hi ha vora el \VUgras{número} de la casa.
Anava a dir-ho en veu alta quan es va sentir el so alegre d'una banda militar. Ens
vam posar a córrer cap al regiment, sense pensar que passaria per davant de les
botigues del \VUgras{barreter} \textsuperscript{(39)} i del \VUgras{guanter}
\textsuperscript{(40)}, i que hauríem estat igual de ben situats per veure'l
desfilar quedant-nos davant de l'aparador de l'\VUgras{òptic}
\textsuperscript{(38)}, ple d'\VUgras{ulleres de pinça, ulleres} i
\VUgras{binocles}.

%%K t14-tit-6 sub
\VUpk{10.2pt}{40}{La desfilada.}
\VUsaut{3.0mm}

%%K t14-14-1 p
14. --- Al capdavant del \VUgras{regiment} \textsuperscript{(41)} marxava el
\VUgras{tambor major} \textsuperscript{(42)}, seguit dels \VUgras{tambors}
\textsuperscript{(43)}, els \VUgras{cornetes} \textsuperscript{(44)} i tota la
\VUgras{banda}. El \VUgras{general} \textsuperscript{(45)}, que era a la plaça amb
el seu ajudant de camp, s'havia aturat vora el \VUgras{brollador}
\textsuperscript{(49)} de la \VUgras{font} \textsuperscript{(48)} per veure la
\VUgras{desfilada}.

%%K t14-14-2 p
\VUgras{Els oficials de l'estat major} \textsuperscript{(46)}, el \VUgras{coronel} i
el \VUgras{tinent coronel} el saludaven amb l'espasa. \VUgras{Els comandants}
anaven al capdavant dels seus \VUgras{batallons} \textsuperscript{(47)} i cada
\VUgras{capità} manava la seva \VUgras{companyia}, mentre els \VUgras{tinents},
els \VUgras{alferes} i els \VUgras{sotsoficials} ocupaven els seus llocs
habituals vora els seus homes.

%%K t14-15-1 p
15. --- Molts transeünts s'havien reunit a la \VUgras{cruïlla}
\textsuperscript{(50)}; els comerciants --- el \VUgras{paraigüer}
\textsuperscript{(51)}, el \VUgras{tintorer} \textsuperscript{(53)}, l'\VUgras{armer}
\textsuperscript{(54)} --- van sortir a veure els \VUgras{soldats}. Tots els
vehicles --- \VUgras{cotxes de punt} \textsuperscript{(52)}, \VUgras{cotxes
particulars} \textsuperscript{(57)} i també la \VUgras{cuba de reg}
\textsuperscript{(56)} --- es van arrambar a la vorera.

%%K t14-tit-7 sub
\VUpk{10.2pt}{40}{Escenes del carrer.}
\VUsaut{3.0mm}

%%K t14-15-2 p
Un \VUgras{viatjant de comerç} \textsuperscript{(55)} que portava a la mà els seus
\VUgras{maletins de mostres} va creuar el carrer a bon pas, i les persones
assegudes a l'\VUgras{imperial} \textsuperscript{(59)} de l'\VUgras{òmnibus}
\textsuperscript{(58)} es van girar per gaudir de l'espectacle. Un pobre
\VUgras{cantant de carrer} \textsuperscript{(60)} amb el seu \VUgras{orguenet}
\textsuperscript{(61)} va haver d'interrompre la seva \VUgras{cançó}, car el
soroll dels tambors ofegava la seva veu.

%%K t14-16-1 p
16. --- Quan el regiment va arribar davant de la \VUgras{botiga de teixits}
\textsuperscript{(64)}, les \VUgras{cosidores} \textsuperscript{(63)} i els
\VUgras{sastres} \textsuperscript{(62)} van deixar les seves \VUgras{màquines de
cosir} i la seva labor per abocar-se a la finestra. \VUgras{El comerciant}
\textsuperscript{(65)}, que acabava de posar \VUgras{vestits}
\textsuperscript{(66)} nous al seu \VUgras{aparador}, va haver de retirar les
peces de \VUgras{drap} \textsuperscript{(67)} i de \VUgras{tela} exposades a la
vorera, vora la \VUgras{boca de claveguera} \textsuperscript{(68)}, per por que la
gent no les tombés; fins i tot un vell \VUgras{marxant}
\textsuperscript{(69)}, a qui una portera regatejava unes mitges, va haver de
ficar-se en un portal per acabar la seva venda.

%%K t14-16-2 p
Vam acompanyar el regiment fins a la caserna \VUgras{i}, molt contents del nostre
dia, vam tornar a casa a l'hora de sopar.

%%K t14-tit-8 sub
\VUpk{10.2pt}{40}{Oficis i professions diversos.}
\VUsaut{3.0mm}

%%K t14-17-1 p
17. --- L'endemà el meu pare es va oferir a portar-nos a visitar la impremta del
diari local. Ho vam acceptar amb entusiasme.

%%K t14-17-2 p
L'\VUgras{impressor} és també \VUgras{llibreter} \textsuperscript{(70)},
\VUgras{editor}, \VUgras{paperer} i \VUgras{enquadernador}. Ha instal·lat al
primer pis la \VUgras{redacció} \textsuperscript{(71)} del diari, i també el
\VUgras{taller de gravat}, on treballen els \VUgras{litògrafs}
\textsuperscript{(72)} i els \VUgras{gravadors en coure} \textsuperscript{(73)}, i
finalment la sala de composició, on són els \VUgras{caixistes}
\textsuperscript{(74)}. La \VUgras{sala de màquines} \textsuperscript{(75)}
pròpiament dita, les \VUgras{premses} i les diverses màquines, són a la planta
baixa.

%%K t14-tit-9 sub
\VUpk{10.2pt}{40}{Joan fa moltíssimes preguntes.}
\VUsaut{3.0mm}

%%K t14-18-1 p
18. --- En sortir de la impremta, Joan es va aturar, molt intrigat, davant d'una
caixeta de vidre muntada damunt un pal davant de la \VUgras{merceria}
\textsuperscript{(77)}, i va preguntar per a què servia. El meu pare li va
explicar que era un \VUgras{avisador d'incendis} \textsuperscript{(76)}. En efecte,
tot a la ciutat era una meravella per al meu amic, des de la senzilla
\VUgras{font de beure} \textsuperscript{(78)} i el lleuger \VUgras{tricicle de
repartiment} \textsuperscript{(80)} que muntava un \VUgras{grum}
\textsuperscript{(79)} de lliurea, fins al \VUgras{furgó de correus}
\textsuperscript{(81)}.

%%K t14-19-1 p
19. --- Mentre el meu pare ens deixava un moment per parlar amb el
\VUgras{recaptador de contribucions} \textsuperscript{(83)}, que s'havia aturat a
conversar amb un \VUgras{advocat} \textsuperscript{(82)} i un \VUgras{notari}
\textsuperscript{(84)}, ens vam entretenir seguint el moviment del carrer. Passava
davant nostre la gent més variada: un \VUgras{aprenent de pastisser}
\textsuperscript{(85)}, que portava en un cistell un esplèndid \VUgras{pastís},
venia darrere d'un \VUgras{drapaire} \textsuperscript{(86)}; un \VUgras{fanaler}
\textsuperscript{(87)} parlava amb un \VUgras{gasista} \textsuperscript{(88)}; una
\VUgras{monja} \textsuperscript{(89)} que portava de la mà una orfeneta tornava al
seu convent.

%%K t14-tit-10 sub
\VUpk{10.2pt}{40}{De tornada a casa.}
\VUsaut{3.0mm}

%%K t14-20-1 p
20. --- Just quan el meu pare tornava amb nosaltres, Joan va deixar anar una
riallada en veure les cares emmascarades i la fatxa estranya de dos
\VUgras{escura-xemeneies} \textsuperscript{(91)} negres de sutge.

%%K t14-21-1 p
21. --- Jo volia comprar una planta per a la meva mare a la \VUgras{florista}
\textsuperscript{(92)}, però el meu pare va observar que pesaria molt per
portar-la i que fóra millor comprar unes \VUgras{taronges}. Ens vam acostar a la
\VUgras{venedora de la parada} \textsuperscript{(94)}, i quan la \VUgras{institutriu}
\textsuperscript{(93)}, que n'estava escollint algunes per al seu infant, va acabar
la seva compra, ens van servir.

%%K t14-22-1 p
22. --- De tornada a casa ens vam trobar amb diversos coneguts: una vella amiga
de la meva família, recolzada al braç de la seva \VUgras{senyoreta de companyia}
\textsuperscript{(95)}, l'\VUgras{enginyer} \textsuperscript{(96)} de camins i
ponts, i una de les meves cosines amb la seva \VUgras{mainadera}
\textsuperscript{(98)}.

%%K t14-22-2 p
Després vam passar vora la \VUgras{modista} \textsuperscript{(97)} de la meva mare
i dos \VUgras{frares} \textsuperscript{(99)} forasters, que es van acostar a
preguntar al meu pare el camí de l'estació.
\VUsaut{6.0mm}
\VUfilet{22mm}
